\section{Limitations}\label{sec:limitations}

\paragraph{Human judgement}
The listening evidence is one native listener, the author, rating six sentences per system blind. It agrees with two
classifiers on the decision that matters, but it is not a listening test. A blind test with multiple native Egyptian
listeners rating naturalness, Egyptian authenticity, similarity and correctness, with real speech as a hidden
anchor, is the next step \pending{and will be reported in a revision}.

\paragraph{Reference voices}
FLEURS speakers read prepared text, mostly in an MSA register; only some of their clips are labelled Egyptian by the
classifiers. Cloning colloquial Egyptian from such references is harder than from everyday speech, and our numbers
may understate what all systems do with a colloquial reference. Six Egyptian voices, three per sentence, also make
the intervals wide.

\paragraph{Classifiers as dialect measures}
Both classifiers label real colloquial Egyptian speech as Egyptian most of the time, but they are trained on
broadcast and online speech, not on synthetic speech, and they disagree on many single utterances. We therefore use
system-level shares, calibrate the bar on real speech, and report both classifiers.

\paragraph{Scope}
\BA adds no training; dialect features that \vox cannot produce, \BA cannot add. Results are for Egyptian Arabic,
English and French with the systems available in November 2026. The text normaliser covers Cairene usage; other
Egyptian varieties (for example Upper Egyptian) are not covered.

\paragraph{Misuse}
Despite the consent step and the watermark, a determined user can clone a voice without permission. We release
because the same technology is already public in less careful forms, and because an open, measured Egyptian system
helps people build, audit and detect such systems; this is a judgement, not a guarantee.
